\newpage
\begin{center}
{\fontsize{15pt}{18pt}\selectfont \textbf{BUILDING A HYBRID RECOMMENDATION MODEL COMBINING MATRIX FACTORIZATION AND DEEP NEURAL NETWORKS}\par}
\vspace{0.5cm}
{\fontsize{16pt}{19.2pt}\selectfont \textbf{ABSTRACT}\par}
\end{center}
\vspace{0.2cm}

Matrix Factorization (MF) models customer--item interactions with the inner product of two latent vectors \parencite{koren2009matrix}: a linear combination that can only be learned for items somebody has already bought. On fashion data -- extremely sparse, with new products launched all the time -- both are limitations. This project builds \textbf{NeuMF-F}, a hybrid model that keeps the two-branch structure of Neural Matrix Factorization \parencite{he2017ncf} -- a linear Generalized Matrix Factorization (GMF) branch and a non-linear Multi-Layer Perceptron (MLP) branch fused at the prediction layer -- but represents every customer and item by an ID embedding plus a projection of features: item attributes and text description, customer attributes, sales and time. During training, item IDs are randomly dropped so that the model learns to score items that have never been sold.

Experiments use the H\&M Personalized Fashion Recommendations dataset \parencite{hm2022kaggle} (31.8 million transactions). Two disjoint, reproducible customer samples are drawn by hashing: development sample A (21,599 customers) is used only for tuning on its validation set, and holdout sample B (21,030 customers) is used once, for the final evaluation. Data are split by time; the $k$-core filter uses only data before the test cut-off; the candidate set includes new items (\VData{Btest}{newpct}\% of the test targets of sample B were never sold before the cut-off); all features are computed point-in-time, without future information. Fourteen models -- NeuMF-F, its GMF-F and MLP-F branches, the late-fusion LateFusion-F, ID-only NeuMF with its two branches, and seven baselines (Random, Most Popular, recent popularity, content-based, ItemKNN, UserKNN, BPR-MF) -- are tuned with the same budget. The final evaluation follows a pre-registered plan with a locked test set: full ranking, five seeds, per-user Wilcoxon signed-rank tests with bootstrap confidence intervals and Holm correction over a family of ten comparisons.

On the test set of sample B (\VData{Btest}{users} customers), NeuMF-F reaches NDCG@10 = \VRes{NeuMFF}{NDCG10}, compared with \VRes{NeuMF}{NDCG10} for ID-only NeuMF, \VRes{BPR}{NDCG10} for BPR-MF, \VRes{UKNN}{NDCG10} for UserKNN and \VRes{RecPop}{NDCG10} for recent popularity. In the pre-registered family, \VSig{all}{summaryEN}. On new items, NeuMF-F reaches NDCG@10 = \VRes{NeuMFF}{NDCG10new}, whereas every ID-only model scores zero by construction. The earlier development phase (ID-only NeuMF, evaluated on the development sample itself) is reported as development history: NeuMF did not outperform a tuned MF baseline, which is why the project turned to adding information rather than only model complexity. An inference-only demo (FastAPI) compares the models for individual customers.
